### Literal statement (as written)

Let (p(x)\in\mathbb{Q}[x]). Define
[
A:={,p(n)+1/n : n\in\mathbb{N},}.
]
Call (A) **strongly complete** if:

> For every finite set (B), the set of all finite subset-sums of (A\setminus B),
> [
> S(A\setminus B):=\left{\sum_{t\in X} t \ :\ X\subseteq A\setminus B \text{ finite}\right},
> ]
> contains all sufficiently large integers, i.e.
> [
> \forall, B\text{ finite}\ \exists N_B\in\mathbb{N}\ \forall m\in\mathbb{Z}\ (m\ge N_B \implies m\in S(A\setminus B)).
> ]

**Ambiguity/misstatement 1 (dummy variable):** The problem text writes (\sum_{n\in X} n) with (X\subseteq A\setminus B). Since (X) consists of elements of (A) (rationals), the summand must be an element of (X), not an index (n). The minimal correction is to write (\sum_{t\in X} t).

**Ambiguity/misstatement 2 ((\mathbb{N}) and (1/n)):** Since (1/n) appears, we must take (\mathbb{N}={1,2,3,\dots}) (otherwise (n=0) is invalid).

**Ambiguity/misstatement 3 (missing growth/positivity hypothesis on (p)):** As written, (p) is arbitrary in (\mathbb{Q}[x]), including polynomials with negative leading coefficient. Then (p(n)+1/n) can be eventually negative and the claim becomes trivially false (see below).

### Minimal corrected statement consistent with standard conventions

A minimal fix consistent with “complete sequences” conventions is to assume (p(n)+1/n>0) for all sufficiently large (n), e.g. that (p) has **positive leading coefficient** (so (p(n)\to+\infty)):

> If (p\in\mathbb{Q}[x]) has positive leading coefficient, is (A={p(n)+1/n:n\ge1}) strongly complete?

I will **disprove the literal statement** (which quantifies over all (p\in\mathbb{Q}[x]) with no sign restriction).

---
